\centerline{\Huge \bf  Экзаменационная работа по алгебре 8 класс}
\centerline{\Huge \bf курс углубленного изучения математики}
\centerline{\Large \bf 1 вариант}

\begin{flushright}
    \scriptsize{Логика может привести Вас от пункта А к пункту Б,\\ а воображение — куда угодно.\\ Альберт Эйнштейн.}
\end{flushright}

\problem{1}
Упростите выражение
\[ \left( \dfrac{1}{a^2 - x^2} + \dfrac{a+x}{x-a}\right)\cdot\dfrac{a+x}{x^2 - a^2 + x-a} + \dfrac{1-2a}{(a-x)^2}\]
\problem{2}
Между какими целыми числами находится каждый корень уравнения

\noindent
$\frac{2}{9x^2 - 4} - \frac{1}{9x^2 - 6x} + \frac{3x - 4}{9x^2 + 6x} = 0$

\problem{3} 
Решить уравнение\\ 
$x(x+1)(x+2)(x+3)=0$
 
\problem{4}
\textbf{Часть 1.} Постройте график функции $f(x) = \sqrt{x^2+4x+4}-1$ и укажите: а) нули функции; б) область определения; в) область значений функции. 

\noindent
\textbf{Часть 2.} Постройте график функции $g(x) = \dfrac{x+2}{x+2} - \dfrac{2}{x-1}$ и укажите: а) нули функции; б) область определения; в) область значений функции.

\noindent
\textbf{Часть 3.} Укажите число точек пересечения графиков $f(x)$ и $g(x).$

\problem{5}
Найдите последнюю цифру числа: 
$3 \cdot 7^{11} - 13 \cdot 2^9 +8$

\problem{6}
\text{Докажите равенство:} 
\[\left( \frac{6 + 4\sqrt{2}}{\sqrt{2} + \sqrt{6 + 4\sqrt{2}}} + \frac{6 - 4\sqrt{2}}{\sqrt{2} - \sqrt{6 - 4\sqrt{2}}} \right)^2 = 2.\]

\problem{7}
При каком значении параметра $m$ сумма квадратов корней уравнения  $x^2+(m-1)x+m^2-1,5=0$ наибольшая?

\problem{8}
\text{Решить совокупность неравенств:}
$
\left[ 
  \begin{gathered} 
    \left\{ 
      \begin{gathered} 
        \dfrac{2x-1}{2}<\dfrac{x+2}{3}
        \\  
        \dfrac{2x+3}{2}\geq \dfrac{x}{4}+\dfrac{1}{2}
        \\
      \end{gathered} 
    \right.  
    \\ 
    |x-1| \leq \dfrac{3}{4}
  \end{gathered}
  \right.
$

\problem{9}
Из пункта А в пункт В отправился велосипедист. Одновременно навстречу ему из пункта В отправился другой велосипедист. При встрече оказалось, что первый проехал на $6$ км меньше второго. Продолжая движение, первый велосипедист прибыл в пункт В через $2$ ч $24$ мин после встречи, а второй прибыл в пункт А через $1$ ч $40$ мин после встречи. Определите, на каком расстоянии от пункта А произошла встреча?

\newpage

\hspace{3mm}

\problem{10}
Найдите натуральные числа $m$ и $n$, если известно, что из трёх следующих утверждений два истины, а одно --- ложно:

1) $4m+9n = 135;$ \qquad 2) $9m+4n = 135$ \qquad 3) $6m+11n=240.$